\chapter{Modèles génératifs et données synthétiques}\label{ch:ml:generative-models-and-synthetic-data}

Un \omterm{def:ml:generative-models-and-synthetic-data:gan}{réseau antagoniste génératif} entraîné sur vingt ans de rendements quotidiens d'indices produit des fenêtres sur lesquelles une stratégie simple,
entraînée et testée sur des données générées, affiche un ratio de Sharpe de 10.7. Entraînée sur les mêmes données générées et
testée sur les dix années réelles qui suivent, elle affiche 0.5. Le réseau avait inventé une prévisibilité que le marché n'avait pas,
et il n'avait pas appris celle que le marché avait. Les données synthétiques comptent parmi les outils les plus utiles de l'apprentissage automatique en
finance~: elles permettent à un desk de soumettre une stratégie à des krachs survenus une seule fois, de tester un moteur de backtest sur des trajectoires dont la vérité est connue,
et d'entraîner des modèles là où les données sont rares. Elles induisent aussi en erreur de façons qu'aucune liste de contrôle des faits stylisés ne repère. Ce chapitre
construit cinq générateurs, depuis le bootstrap par blocs jusqu'au \omterm{def:ml:generative-models-and-synthetic-data:diffusion}{modèle de diffusion}, et les juge au moyen de quatre tests~; ses données sont
synthétiques, afin que la série réelle possède une structure connue à apprendre.

\section{À quoi sert un modèle génératif}

\begin{definition}[Modèle génératif]\label{def:ml:generative-models-and-synthetic-data:gen}
Un \emph{modèle génératif}\index{modèle génératif} est un modèle de la distribution des données qui peut tirer de nouveaux échantillons
de celle-ci. Pour les données de marché, les échantillons sont des trajectoires (de rendements, d'états du carnet d'ordres, de courbes) ou des fenêtres de ces trajectoires, et le modèle est jugé
selon que les échantillons sont utiles à l'usage visé, et non seulement selon qu'ils paraissent réels.
\end{definition}

La série réelle du chapitre est constituée de trente ans de rendements quotidiens d'indice issus d'un modèle GJR-GARCH à chocs de loi de Student (livre~4,
chapitre~18)~: queues épaisses, regroupement de volatilité, effet de levier, six krachs introduits de 8 à 11 écarts types,
et une part de prévisibilité~: une moyenne conditionnelle égale à 0.1 fois la volatilité courante, dans le sens du rendement des
vingt derniers jours. Les vingt premières années servent à l'entraînement~; les dix dernières servent au test. Chaque générateur produit des fenêtres de 32 jours de bourse. Ses
usages fixent les tests~: un test de résistance exige les queues et le regroupement~; entraîner un prédicteur exige la prévisibilité,
et qu'elle ne soit pas inventée~; la confidentialité ou le partage des données exigent des échantillons qui ne soient pas des copies.

Deux générateurs classiques servent de référence. Un bootstrap par blocs (livre~4, chapitre~13) découpe la série d'entraînement en blocs de 20 jours
puis recolle des blocs tirés au hasard~: il conserve tout à l'intérieur d'un bloc et rien de ce qui relie deux blocs. Un modèle GARCH(1,1)
à chocs de Student, estimé par maximum de vraisemblance ($\alpha = 0.079$, $\beta = 0.886$, 3.6 degrés de liberté), conserve
ce que disent ses équations et rien d'autre~: il n'a ni effet de levier ni momentum.

\section{Autoencodeurs variationnels et réseaux antagonistes}

\begin{definition}[Autoencodeur variationnel]\label{def:ml:generative-models-and-synthetic-data:vae}
Un \emph{autoencodeur variationnel}\index{autoencodeur variationnel} (VAE) est un \omterm{def:ml:cross-sectional-deep-models:ae}{autoencodeur} (chapitre~9) dont l'encodeur produit une
distribution sur un code latent de faible dimension et dont le décodeur produit une distribution sur les données~; il est entraîné à
maximiser une borne inférieure de la vraisemblance, à savoir la log-vraisemblance de reconstruction moins la divergence de Kullback--Leibler
de la distribution de l'encodeur par rapport à une loi a priori normale centrée réduite, et il génère en décodant des codes tirés de cette loi a priori (Kingma et
Welling, 2014).
\end{definition}

\begin{definition}[Réseau antagoniste génératif, effondrement de modes]\label{def:ml:generative-models-and-synthetic-data:gan}
Un \emph{réseau antagoniste génératif}\index{réseau antagoniste génératif} (GAN) entraîne deux réseaux l'un contre l'autre~:
un générateur qui transforme du bruit en échantillons, et un discriminateur qui distingue les échantillons générés des réels~; le générateur est
entraîné à tromper le discriminateur (Goodfellow et coauteurs, 2014). L'\emph{effondrement de modes}\index{effondrement de modes} (mode collapse) est la
défaillance dans laquelle le générateur ne couvre qu'une partie de la distribution des données (quelques motifs typiques, sans les extrêmes)
parce que cela suffit à tromper le discriminateur courant.
\end{definition}

L'\omterm{def:ml:generative-models-and-synthetic-data:vae}{autoencodeur variationnel} du chapitre a un code latent à quatre dimensions et produit une moyenne et une variance pour chacun des 32 jours~; son GAN
transforme un bruit à 16 dimensions à travers deux couches de 128 unités. Les deux sont des \omterm{def:ml:neural-networks-for-noisy-tabular-data:mlp}{perceptrons multicouches} sur des fenêtres standardisées par la
volatilité des rendements d'entraînement, et entraînés pendant 4\,000 pas d'\omterm{def:ml:neural-networks-for-noisy-tabular-data:adam}{Adam} sur un cœur en quelques secondes. Les fenêtres du VAE ont une kurtosis en excès
de 11.1 contre 13.0 pour les fenêtres réelles d'entraînement~; celles du GAN n'atteignent que 2.0~: il produit des semaines calmes ou modérément volatiles
et jamais de krach, signature de l'\omterm{def:ml:generative-models-and-synthetic-data:gan}{effondrement de modes} sur une distribution à queue épaisse. Les Quant GANs (Wiese et coauteurs, 2020) et
le simulateur de marché à petits échantillons de Buehler et coauteurs (2020) sont des conceptions visant ces défaillances.

\section{Modèles de diffusion}

\begin{definition}[Modèle de diffusion]\label{def:ml:generative-models-and-synthetic-data:diffusion}
Un \emph{modèle de diffusion}\index{modèle de diffusion} apprend à inverser un processus qui ajoute progressivement du bruit gaussien aux
données~: un réseau est entraîné à prédire le bruit ajouté à un échantillon à chacune des $T$ étapes, et de nouveaux échantillons sont engendrés en
partant d'un bruit pur et en retirant pas à pas le bruit prédit (Ho, Jain et Abbeel, 2020).
\end{definition}

Le \Cref{lst:ml:gen:diffusion} présente un \omterm{def:ml:generative-models-and-synthetic-data:diffusion}{modèle de diffusion} débruitant complet sur des fenêtres~: 100 niveaux de bruit, un réseau à deux couches qui
prend la fenêtre bruitée et le pas, ainsi que la chaîne inverse. C'est le plus stable des trois réseaux à entraîner, et sur ces données
le plus faible dans les queues (kurtosis de 4.0)~: le bruit gaussien dont il part et les quelques milliers de pas d'entraînement
laissent les extrêmes sous-représentés.

\section{Juger les données synthétiques}

\begin{definition}[Test à deux échantillons par classifieur, entraîner sur le synthétique, tester sur le réel]\label{def:ml:generative-models-and-synthetic-data:tests}
Le \emph{test à deux échantillons par classifieur}\index{test à deux échantillons par classifieur} entraîne un classifieur à distinguer les échantillons générés des
réels et en rapporte l'exactitude sur des échantillons mis de côté~: 50~\% signifie que le classifieur ne peut pas les distinguer (Lopez-Paz et
Oquab, 2017). \emph{Entraîner sur le synthétique, tester sur le réel}\index{entraîner sur le synthétique, tester sur le réel} (TSTR) entraîne un modèle aval
sur des données générées et l'évalue sur des données réelles, en le comparant au même modèle entraîné sur des données réelles (Esteban, Hyland et
Rätsch, 2017).
\end{definition}

Quatre tests jugent les cinq générateurs (\cref{tab:ml:gen:judge}). La \emph{fiche de notation} vérifie six faits stylisés (Cont,
2001) à l'intérieur des fenêtres~: queues épaisses (excès de kurtosis supérieur à 5), absence d'autocorrélation linéaire, regroupement de volatilité
dans les rendements en valeur absolue aux retards 1 et 10, l'effet de levier (corrélation négative d'un rendement avec le rendement absolu suivant), asymétrie gains--pertes
(skewness négative) et gaussianité par agrégation (kurtosis des sommes sur 20 jours inférieure à la moitié de celle des rendements quotidiens). Les
fenêtres réelles d'entraînement réussissent les six. Le \emph{test à deux échantillons} utilise un classifieur par \omterm{def:ml:trees-and-boosting:boosting}{gradient boosting} sur 13 variables de fenêtre~; les
fenêtres réelles se chevauchent, donc ses moitiés sont découpées par année, faute de quoi le classifieur reconnaît une fenêtre de test à ses voisines dans la
moitié d'entraînement (une séparation aléatoire donne 0.83 aux années paires d'entraînement contre les impaires). \emph{TSTR} ajuste une
prévision linéaire du dernier jour d'une fenêtre à partir des variables de ses 31 premiers jours, sur 20\,000 fenêtres générées, et trade le signe de la
prévision sur les dix années réelles de test~; elle est moyennée sur cinq échantillons générés. \emph{Entraîner sur le synthétique,
tester sur le synthétique} (TSTS) fait de même avec un second échantillon généré comme \omterm{def:ml:why-financial-machine-learning-is-different:sets}{ensemble de test}. La dernière colonne est le
\emph{contrôle de \omterm{def:ml:large-language-models-in-finance:memo}{mémorisation}} (\cref{lst:ml:gen:tests})~: la part des fenêtres générées plus proches d'une fenêtre d'entraînement que
ne le sont 95~\% des fenêtres réelles de test.

\begin{table}[htbp]
\centering\small
\setlength{\tabcolsep}{4pt}%
\begin{tabular}{@{}>{\raggedright\arraybackslash}p{3.6cm}>{\centering\arraybackslash}p{1.6cm}>{\centering\arraybackslash}p{1.6cm}>{\centering\arraybackslash}p{1.6cm}>{\centering\arraybackslash}p{1.6cm}>{\centering\arraybackslash}p{1.6cm}>{\centering\arraybackslash}p{1.6cm}@{}}
\toprule
& faits & deux échantillons & \multicolumn{2}{c}{TSTR} & TSTS & plus proche que\\
\cmidrule(lr){4-5}
générateur & réussis & exactitude & IC & Sharpe & Sharpe & réel (\%)\\
\midrule
réel (entraîner sur le réel, tester sur le réel) & 6 & 0.545 & 0.078 & 0.86 & & 5.0\\
bootstrap par blocs & 6 & 0.491 & 0.062 & 0.71 & 1.19 & 14.5\\
GARCH-t & 4 & 0.539 & $-0.019$ & $-0.10$ & $-0.09$ & 5.7\\
VAE & 6 & 0.619 & 0.040 & 0.76 & 0.79 & 5.1\\
GAN & 3 & 0.749 & 0.052 & 0.52 & 10.68 & 11.2\\
diffusion & 4 & 0.693 & 0.030 & 0.23 & 1.45 & 1.4\\
\bottomrule
\end{tabular}
\caption{Cinq générateurs jugés de quatre façons. Exactitude du test à deux échantillons sur la ligne réelle~: les années paires
d'entraînement contre les années impaires. Données~: \texttt{ml\_gen.judge}.}\label{tab:ml:gen:judge}
\end{table}

\begin{omfigure}
\begin{tikzpicture}
\begin{axis}[width=11.5cm,height=6cm,xlabel={retard (jours)},ylabel={autocorrélation de $|r_t|$},tick label style={font=\small},
  label style={font=\small},xmin=1,xmax=20,ymin=0,ymax=0.32,grid=major,grid style={gray!25},legend columns=3,
  legend style={font=\footnotesize,at={(0.5,-0.25)},anchor=north,draw=none,fill=none,
  /tikz/every even column/.append style={column sep=6pt}},table/col sep=comma]
\addplot[black,ultra thick] table[x=lag,y=real]{figdata/ml/16-generative-models-and-synthetic-data/absacf.csv};
\addplot[omDef,thick,mark=*,mark size=1.3pt] table[x=lag,y=boot]{figdata/ml/16-generative-models-and-synthetic-data/absacf.csv};
\addplot[gray,thick,mark=triangle*,mark size=1.3pt] table[x=lag,y=garch]{figdata/ml/16-generative-models-and-synthetic-data/absacf.csv};
\addplot[omProp,thick,mark=square*,mark size=1.3pt] table[x=lag,y=vae]{figdata/ml/16-generative-models-and-synthetic-data/absacf.csv};
\addplot[omThm,thick,mark=diamond*,mark size=1.5pt] table[x=lag,y=gan]{figdata/ml/16-generative-models-and-synthetic-data/absacf.csv};
\addplot[omProp!50,thick,dashed] table[x=lag,y=diff]{figdata/ml/16-generative-models-and-synthetic-data/absacf.csv};
\legend{réel,bootstrap par blocs,GARCH-t,VAE,GAN,diffusion}
\end{axis}
\end{tikzpicture}

\omcaption{Regroupement de volatilité~: autocorrélation des rendements quotidiens absolus dans des fenêtres de 32~jours, fenêtres réelles d'entraînement
contre 4\,000 fenêtres issues de chaque générateur. Données~: \texttt{ml\_gen.abs\_acf}.}\label{fig:ml:gen:acf}
\end{omfigure}

Le tableau rassemble les enseignements du chapitre. Le scorecard est le test le plus faible~: le VAE passe les six faits stylisés et se laisse distinguer
des fenêtres réelles dans 62~\% des cas~; le bootstrap par blocs passe les six faits et est indiscernable (0.491) parce que ses fenêtres
sont des morceaux de fenêtres réelles. Le test à deux échantillons voit ce que le scorecard rate, mais pas ce qui compte pour une stratégie. Seul le TSTR
mesure cela~: le modèle réel rapporte 0.86~; les fenêtres du bootstrap conservent l'essentiel du momentum à l'intérieur de leurs blocs de 20~jours
(0.71)~; le VAE en a appris une version bruitée (0.76, avec un écart type de 0.28 d'un échantillon à l'autre)~; GARCH, qui ne peut pas la
représenter, n'apprend rien ($-0.10$, écart type de 0.65~: pur bruit)~; le \omterm{def:ml:generative-models-and-synthetic-data:diffusion}{modèle de diffusion} apprend peu (0.23).

\section{Les façons dont les données synthétiques induisent en erreur}

La \Cref{fig:ml:gen:sharpe} est le résultat de l'accroche. Sur les fenêtres propres au GAN, une stratégie obtient un ratio de Sharpe de 10.7~: le
générateur produit une autocorrélation linéaire bien plus forte (0.099 au retard 1, contre 0.078 dans les fenêtres réelles) et des
profils persistants de volatilité qu'un modèle linéaire peut exploiter, un artefact d'un réseau qui s'est effondré sur quelques formes lisses.
Testée sur des données réelles, la même stratégie rapporte 0.52. Tous les générateurs, sauf GARCH, présentent un écart entre TSTS et TSTR, et
rien à l'intérieur du monde synthétique ne le révèle.

\begin{omfigure}
\begin{tikzpicture}
\begin{axis}[width=11cm,height=6cm,ybar,bar width=10pt,ymin=-1,ymax=11.5,ylabel={ratio de Sharpe (annualisé)},area legend,
  xtick={0,1,2,3,4},xticklabels={{bootstrap\\par blocs},GARCH-t,VAE,GAN,diffusion},xticklabel style={align=center},tick label style={font=\small},
  label style={font=\small},xmin=-0.5,xmax=4.5,grid=major,grid style={gray!25},legend columns=3,
  legend style={font=\footnotesize,at={(0.5,-0.2)},anchor=north,draw=none,fill=none},table/col sep=comma]
\addplot[fill=omThm,draw=none,bar shift=-5.5pt] table[x=i,y=tsts]{figdata/ml/16-generative-models-and-synthetic-data/sharpe.csv};
\addplot[fill=omDef,draw=none,bar shift=5.5pt] table[x=i,y=tstr]{figdata/ml/16-generative-models-and-synthetic-data/sharpe.csv};
\addplot[black!70,dashed,thick,sharp plot,forget plot] table[x=x,y=y]{figdata/ml/16-generative-models-and-synthetic-data/trtr.csv};
\legend{testé sur les fenêtres générées,testé sur les fenêtres réelles}
\end{axis}
\end{tikzpicture}

\omcaption{La même stratégie, entraînée sur les fenêtres de chaque générateur, testée sur un second échantillon généré et sur les dix
années réelles de test~; en pointillés, entraînée et testée sur des données réelles. Données~: \texttt{ml\_gen.judge}.}\label{fig:ml:gen:sharpe}
\end{omfigure}

La colonne de \omterm{def:ml:large-language-models-in-finance:memo}{mémorisation} montre les deux autres échecs. Les fenêtres du bootstrap par blocs sont plus proches d'une fenêtre d'entraînement que
ne le sont les fenêtres réelles de test, trois fois plus souvent que par hasard (14.5~\% contre 5~\%)~: ce sont des copies par construction, ce qui
est acceptable pour les tests de résistance et inutile là où les données synthétiques doivent protéger le propriétaire des données réelles. Le 11.2~\% du GAN a une
cause différente~: toutes ses fenêtres sont calmes, et les fenêtres calmes ont des voisines proches parmi les nombreuses fenêtres calmes d'entraînement.
La distance au plus proche voisin signale aussi bien la copie que l'effondrement, et doit être lue avec les queues de distribution.
La \omterm{def:ml:large-language-models-in-finance:memo}{mémorisation} du chapitre~14 s'applique aux générateurs comme aux \omterm{def:ml:large-language-models-in-finance:lm}{modèles de langage}~: un grand générateur entraîné assez longtemps sur un petit historique de marché
le reproduit.

\begin{method}[Utiliser des données de marché synthétiques]\label{met:ml:gen:method}
\begin{enumerate}
\item Nommer la finalité (résistance, entraînement, test d'un moteur de backtest, partage), et choisir les tests dont elle a besoin.
\item Partir des générateurs classiques (bootstrap, GARCH, un simulateur aux mécanismes connus) et exiger qu'un réseau les
batte sur ces tests.
\item Juger par le TSTR, face à l'entraînement sur le réel, jamais par les résultats à l'intérieur du monde synthétique~; ajouter le test à deux échantillons avec une
découpe par blocs et le contrôle de \omterm{def:ml:large-language-models-in-finance:memo}{mémorisation}.
\item Traiter tout edge trouvé seulement dans les données synthétiques comme un artefact tant qu'il n'est pas retrouvé dans les données réelles.
\end{enumerate}
\end{method}

\section{Tutoriel~: des trajectoires qui ne se sont jamais produites}\label{tut:ml:generative-models-and-synthetic-data:tut}

\begin{tutorial}
\textbf{Objectif.} Ajuster cinq générateurs sur les années d'entraînement, et les juger au moyen du scorecard, du test à deux
échantillons, du TSTR, du TSTS et de la distance au plus proche voisin. \textbf{État final~:} le \cref{tab:ml:gen:judge}, les \cref{fig:ml:gen:acf,fig:ml:gen:sharpe}.
\begin{tutsteps}
\item \textbf{Un \omterm{def:ml:generative-models-and-synthetic-data:gan}{réseau antagoniste génératif}.}
\omcode{code/firm/genmkt/firm_genmkt.py}{166}{195}{Un générateur et un discriminateur entraînés l'un contre l'autre.}\label{lst:ml:gen:gan}
\item \textbf{Un \omterm{def:ml:generative-models-and-synthetic-data:diffusion}{modèle de diffusion} par débruitage.}
\omcode{code/firm/genmkt/firm_genmkt.py}{198}{235}{Le bruit, un débruiteur et la chaîne inverse.}\label{lst:ml:gen:diffusion}
\item \textbf{\omterm{def:ml:generative-models-and-synthetic-data:tests}{Entraîner sur le synthétique, tester sur le réel}.}
\omcode{code/firm/genmkt/firm_genmkt.py}{308}{319}{Le score TSTR~: IC de rang et ratio de Sharpe sur les fenêtres réelles.}\label{lst:ml:gen:tests}
\item \textbf{Exécuter} \texttt{ml\_gen.judge()} (environ une minute sur un cœur) et \texttt{fig\_gen.py}.
\end{tutsteps}
\textbf{Pour aller plus loin.} Conditionner les générateurs au régime de volatilité~; entraîner le GAN plus longtemps et surveiller son TSTS
ratio de Sharpe et sa précision au test à deux échantillons~; remplacer les blocs de 20~jours du bootstrap par blocs par des blocs de 5~jours et observer
la disparition du momentum dans son TSTR.
\end{tutorial}

\section{Construction~: générateurs et leurs juges}\label{bld:ml:generative-models-and-synthetic-data:genmkt}

\begin{build}
\bfield{Rôle} Données de marché synthétiques dont la finalité est énoncée, jugées par des tests qui y correspondent.
\bfield{Interface} \texttt{real\_index(n\_days, seed)}, \texttt{windows}~; \texttt{block\_bootstrap}, \texttt{fit\_garch},
\texttt{garch\_windows}~; \texttt{VAE}, \texttt{GAN}, \texttt{Diffusion} avec \texttt{fit(W, steps, seed)} et
\texttt{sample(n, seed)}~; \texttt{facts}, \texttt{scorecard}, \texttt{c2st(A, B, groups\_a, groups\_b, seed)},
\texttt{tstr(W\_train, W\_test)}, \texttt{nn\_distance}.
\bfield{Règles} Chaque générateur est ajusté sur la seule période d'entraînement~; les fenêtres réelles qui se chevauchent sont découpées dans le temps~; le TSTR est
moyenné sur les échantillons~; les résultats à l'intérieur du monde synthétique ne sont jamais rapportés seuls.
\bfield{Tests d'acceptation} \texttt{code/firm/genmkt/tests/}~: l'ajustement GARCH retrouve les paramètres injectés~; le bootstrap recopie
des blocs~; le test à deux échantillons est proche de 0.5 sur deux échantillons d'une même distribution et proche de 1 sur deux distributions différentes~; le TSTR retrouve un
prédicteur injecté~; les réseaux s'entraînent de façon déterministe et réduisent leurs pertes~; la distance au plus proche voisin est nulle pour
les copies.
\bfield{Pour aller plus loin} Génération conditionnelle selon les régimes~; générateurs fondés sur la signature~; le marché à base d'agents du livre~10 (chapitre~27)
comme générateur doté de mécanismes, jugé par les mêmes tests.
\end{build}

\begin{omsources}
\item D.~P.~Kingma et M.~Welling, ``Auto-encoding variational Bayes'', arXiv:1312.6114, 2014.
\item I.~Goodfellow et coauteurs, ``Generative adversarial nets'', arXiv:1406.2661, 2014.
\item J.~Ho, A.~Jain et P.~Abbeel, ``Denoising diffusion probabilistic models'', arXiv:2006.11239, 2020.
\item M.~Wiese, R.~Knobloch, R.~Korn et P.~Kretschmer, ``Quant GANs: deep generation of financial time series'',
\emph{Quantitative Finance} 20(9), 2020.
\item H.~Buehler, B.~Horvath, T.~Lyons, I.~Perez Arribas et B.~Wood, ``A data-driven market simulator for small data
environments'', arXiv:2006.14498, 2020.
\item D.~Lopez-Paz et M.~Oquab, ``Revisiting classifier two-sample tests'', arXiv:1610.06545, 2017.
\item R.~Cont, ``Empirical properties of asset returns: stylized facts and statistical issues'', \emph{Quantitative Finance}
1(2), 2001.
\item C.~Esteban, S.~L.~Hyland et G.~Rätsch, ``Real-valued (medical) time series generation with recurrent conditional
GANs'', arXiv:1706.02633, 2017.
\end{omsources}

\section{Exercices}

\begin{exercise}[$\star$]\label{exo:ml:generative-models-and-synthetic-data:1}
Un test à deux échantillons obtient un score de 0.62 sur 4\,000 fenêtres générées contre 4\,000 fenêtres réelles (2\,000 de chaque réservées au test). Dire si l'écart
par rapport à 0.5 est significatif, et ce que signifierait un score de 0.51.
\end{exercise}

\begin{exercise}[$\star$]\label{exo:ml:generative-models-and-synthetic-data:2}
Indiquer lesquels des six faits stylisés un modèle GARCH(1,1) à chocs de Student symétriques ne peut jamais produire, et expliquer pourquoi.
\end{exercise}

\begin{exercise}[$\star$]\label{exo:ml:generative-models-and-synthetic-data:3}
Sachant que le modèle GARCH ajusté vérifie $\alpha + \beta = 0.964$, calculer la demi-vie d'un choc de volatilité, en jours.
\end{exercise}

\begin{exercise}[$\star\star$]\label{exo:ml:generative-models-and-synthetic-data:4}
Expliquer pourquoi le bootstrap par blocs conserve en partie l'effet de momentum, et indiquer quelle longueur de bloc le détruirait.
\end{exercise}

\begin{exercise}[$\star\star$]\label{exo:ml:generative-models-and-synthetic-data:5}
Expliquer pourquoi un découpage aléatoire a conduit le test à deux échantillons à obtenir 0.83 en comparant des fenêtres réelles à des fenêtres réelles.
\end{exercise}

\begin{exercise}[$\star\star$]\label{exo:ml:generative-models-and-synthetic-data:6}
\emph{Trouver la faille.} «~Les trajectoires de notre GAN passent tous nos tests de faits stylisés, et une stratégie entraînée sur un million d'entre elles affiche un
ratio de Sharpe de 3 hors échantillon sur un autre million, elle est donc robuste.~»
\end{exercise}

\begin{exercise}[$\star\star\star$]\label{exo:ml:generative-models-and-synthetic-data:7}
\emph{Programmation.} Entraîner le GAN pendant 1\,000 puis 8\,000 pas. Indiquer, à chaque fois, son kurtosis en excès, son exactitude à deux échantillons et son
ratio de Sharpe TSTS. Un entraînement plus long supprime-t-il l'\omterm{def:ml:generative-models-and-synthetic-data:gan}{effondrement de modes}~?
\end{exercise}

\begin{exercise}[$\star\star\star$]\label{exo:ml:generative-models-and-synthetic-data:8}
Établir l'objectif de l'\omterm{def:ml:generative-models-and-synthetic-data:vae}{autoencodeur variationnel}~: montrer que $\log p(x)\ge\E_{q(z\mid x)}[\log p(x\mid z)] - \mathrm{KL}(q(z\mid x)\,\|\,p(z))$, et
écrire le terme KL pour une gaussienne diagonale $q$ et une loi a priori normale centrée réduite.
\end{exercise}

\section{Problème~: Trajectoires qui ne se sont jamais produites}

\begin{problem}[{Problème du week-end --- la forme ne garantit pas la prévisibilité}]\label{pb:ml:generative-models-and-synthetic-data:1}
Les trente ans de rendements du chapitre, cinq générateurs et quatre tests.

\medskip\noindent\textbf{Partie I --- Les données et les générateurs.}
\begin{enumerate}
\item Quelle structure présente la série réelle, et quelles parties chaque générateur peut-il représenter~?
\item Quels sont les paramètres GARCH ajustés, et que laisse échapper le modèle~?
\item À quoi ressemblent les réseaux VAE et GAN, et combien de temps dure leur entraînement~?
\item Quel est le processus de génération du \omterm{def:ml:generative-models-and-synthetic-data:diffusion}{modèle de diffusion}~?
\end{enumerate}

\medskip\noindent\textbf{Partie II --- Évaluation.}
\begin{enumerate}[resume]
\item Combien de faits stylisés chaque générateur passe-t-il~?
\item Quelles sont les exactitudes à deux échantillons, et pourquoi faut-il découper les fenêtres réelles dans le temps~?
\item Quels sont les coefficients d'information TSTR et les ratios de Sharpe, comparés à ceux d'un entraînement sur le réel~?
\item Que montre le contrôle de \omterm{def:ml:large-language-models-in-finance:memo}{mémorisation}~?
\end{enumerate}

\medskip\noindent\textbf{Partie III --- Pièges.}
\begin{enumerate}[resume]
\item Quel ratio de Sharpe donne l'univers du GAN, et d'où vient-il~?
\item Pourquoi le score TSTR du GARCH est-il bruité, et quelle est l'ampleur de ce bruit~?
\item Quel générateur utilise-t-on pour un test de résistance, et lequel pour enrichir un \omterm{def:ml:why-financial-machine-learning-is-different:sets}{ensemble d'entraînement}~?
\item Pourquoi le bootstrap ne convient-il pas au partage de données en dehors de l'entreprise~?
\end{enumerate}

\medskip\noindent\textbf{Partie IV --- Le verdict.}
\begin{enumerate}[resume]
\item Énoncer le \emph{résultat nommé}~: pour chaque générateur, le nombre de tests réussis sur la grille de faits stylisés et l'exactitude à deux échantillons, ainsi que l'écart TSTR en
ratio de Sharpe.
\item Qu'est-ce qui permettrait de conclure qu'un générateur a appris une vraie prévisibilité~?
\item Comment tester un outil de backtest à l'aide d'un générateur~?
\item Quand un simulateur mécaniste est-il préférable à un générateur appris~?
\item Comment conditionner un générateur à un régime, et quel nouveau risque cela introduit-il~?
\item Que faudrait-il consigner sur un jeu de données synthétique utilisé en recherche~?
\item Qu'implique pour les tests de résistance la faiblesse du \omterm{def:ml:generative-models-and-synthetic-data:diffusion}{modèle de diffusion} dans les queues de distribution~?
\item En une phrase~: qu'apprend un \omterm{def:ml:generative-models-and-synthetic-data:gen}{modèle génératif} à partir de vingt ans de rendements~?
\end{enumerate}
\end{problem}

\section{Questions d'entretien}

\begin{interviewq}[$\star$ \iqroles{chercheur, ingénieur ML}]\label{iq:ml:generative-models-and-synthetic-data:1}
Citez les principaux faits stylisés des rendements quotidiens. Lesquels un modèle GARCH reproduit-il~?
\end{interviewq}

\begin{interviewq}[$\star\star$ \iqroles{ingénieur ML}]\label{iq:ml:generative-models-and-synthetic-data:2}
Expliquez la différence entre un VAE, un GAN et un \omterm{def:ml:generative-models-and-synthetic-data:diffusion}{modèle de diffusion}, en un paragraphe chacun, et donnez un échec typique de chacun.
\end{interviewq}

\begin{interviewq}[$\star\star$ \iqroles{chercheur}]\label{iq:ml:generative-models-and-synthetic-data:3}
Comment évalueriez-vous un générateur de données de marché synthétiques~?
\end{interviewq}

\begin{interviewq}[$\star\star$ \iqroles{chercheur, trader}]\label{iq:ml:generative-models-and-synthetic-data:4}
Un collègue enrichit l'\omterm{def:ml:why-financial-machine-learning-is-different:sets}{ensemble d'entraînement} avec des trajectoires du GAN, et le backtest s'améliore. Que vérifiez-vous~?
\end{interviewq}

\begin{interviewq}[$\star\star$ \iqroles{ingénieur ML}]\label{iq:ml:generative-models-and-synthetic-data:5}
Qu'est-ce que l'\omterm{def:ml:generative-models-and-synthetic-data:gan}{effondrement de modes} (mode collapse), et comment le détecteriez-vous dans les rendements générés~?
\end{interviewq}

\begin{interviewq}[$\star\star\star$ \iqroles{chercheur}]\label{iq:ml:generative-models-and-synthetic-data:6}
Concevez un \omterm{def:ml:generative-models-and-synthetic-data:tests}{test à deux échantillons par classifieur} pour des fenêtres qui se chevauchent dans une série temporelle. Qu'est-ce qui ne va pas avec un découpage aléatoire~?
\end{interviewq}
